% =============================================================================
% Central Insulin Resistance as a Mediator of GLP-1 Receptor Agonist
% Neuroprotection in Alzheimer's and Parkinson's Disease
% Draft v1 -- structure/formatting skeleton with placeholder figures.
% Populate TODO markers as your own analysis/figures become available.
% =============================================================================

\documentclass[11pt, a4paper]{article}

% ---------------------------------------------------------------- PACKAGES --
\usepackage[utf8]{inputenc}
\usepackage[T1]{fontenc}
\usepackage{geometry}
\geometry{margin=1in}
\usepackage{graphicx}
\usepackage{amsmath, amssymb}
\usepackage{booktabs}
\usepackage{caption}
\usepackage{subcaption}
\usepackage{hyperref}
\hypersetup{colorlinks=true, linkcolor=blue, citecolor=blue, urlcolor=blue}
\usepackage[round]{natbib}
\usepackage{setspace}
\usepackage{titlesec}
\usepackage{enumitem}
\usepackage{xcolor}

% Quick visual placeholder for figures not yet attached -- replace each
% \FigurePlaceholder call with a real \includegraphics{figures/...} once you
% have the image. Renders as a labeled empty box so the document still
% compiles cleanly before any real figures exist.
\newcommand{\FigurePlaceholder}[2]{%
  \fbox{\parbox[c][#1][c]{0.92\textwidth}{\centering\color{gray}\itshape #2}}%
}

\onehalfspacing
\titleformat{\section}{\large\bfseries}{\thesection}{1em}{}
\titleformat{\subsection}{\bfseries}{\thesubsection}{1em}{}

% ----------------------------------------------------------------- TITLE ---
\title{%
  Central Insulin Resistance as a Mediator of GLP-1 Receptor Agonist\\
  Neuroprotection in Alzheimer's and Parkinson's Disease:\\
  \large A Systematic Evidence Synthesis and Exploratory Transcriptomic Analysis
}
\author{%
  % TODO: author name(s), affiliation(s), corresponding-author email
  Your Name$^{1}$ \\
  \small $^{1}$Your Institution / Department
}
\date{\today}

% =============================================================================
\begin{document}

\maketitle

% --------------------------------------------------------------- ABSTRACT --
\begin{abstract}
% TODO: finalize last, after Results/Discussion are complete. Draft below
% follows a standard Background/Methods/Results/Conclusion structure.
\noindent\textbf{Background:} Alzheimer's disease (AD) and Parkinson's
disease (PD) have both been linked to impaired brain (central) insulin
signaling. Glucagon-like peptide-1 receptor agonists (GLP-1RAs) cross the
blood--brain barrier and show neuroprotective effects in preclinical and
some clinical studies, with central insulin resistance proposed as the
mediating mechanism. \textbf{Methods:} We conducted a structured literature
synthesis (PubMed/PMC/clinical trial registries, $n=21$ primary sources)
evaluating three candidate hypotheses regarding the role of central insulin
resistance, supplemented by an original differential-expression re-analysis
of a public transcriptomic dataset (GEO accession GSE34451) and a
formal causal mediation-analysis framework. \textbf{Results:} No identified
source, clinical or preclinical, performed a formal statistical mediation
test linking GLP-1RA exposure, a central insulin-resistance marker, and a
neurodegenerative outcome in the same subjects. Mechanistic convergence on
shared insulin-signaling pathways was common across preclinical studies and
reviews (10/21 sources), but insulin-pathway engagement and neuroprotective
phenotype frequently dissociated within individual studies, clinical trial
endpoints were directionally inconsistent between diseases (positive on
primary motor endpoints in 3/3 Parkinson's disease RCTs vs.\ null or mixed
on primary endpoints in the three Alzheimer's disease RCTs reviewed), and
the one direct empirical test of whether metabolic/insulin-resistance status
predicts treatment response (Athauda et al.\ 2019 post-hoc analysis) found
no such association. An original re-analysis of GSE34451 (rat hippocampus,
Control vs.\ type 2 diabetic, $n=3$/group) found that of 125 KEGG
insulin-signaling pathway genes, none survived FDR correction, but five
reached nominal significance (raw $p<0.05$) with mixed direction of effect:
\textit{Mapk8} (JNK1) upregulation in the diabetic group is consistent with
the canonical JNK-mediated insulin-resistance mechanism, whereas
\textit{Insr} and \textit{Gys2} upregulation runs counter to the simplest
version of the hypothesis. \textbf{Conclusions:} Central insulin resistance
is a biologically plausible contributory mechanism in both diseases (rather
than the dominant or subgroup-defining mechanism), but formal mediation
evidence does not yet exist in the literature; closing this gap requires a
dataset or sub-study that measures GLP-1RA exposure, a central
insulin-resistance marker, and a neurodegenerative outcome in the same
subjects.

\vspace{0.5em}
\noindent\textbf{Keywords:} GLP-1 receptor agonists; insulin resistance;
Alzheimer's disease; Parkinson's disease; neuroprotection; mediation analysis
\end{abstract}

% =============================================================================
\section{Introduction}
\label{sec:intro}

\subsection{Background}
Alzheimer's disease (AD) and Parkinson's disease (PD) have both been linked
to impaired brain (central) insulin signaling -- sometimes termed
``type 3 diabetes'' in the AD literature -- contributing to
neuroinflammation, mitochondrial dysfunction, pathological protein
aggregation (amyloid-$\beta$, tau, $\alpha$-synuclein), and synaptic loss
\citep{talbot2014}. GLP-1 receptors are expressed throughout the central
nervous system, including the hippocampus, substantia nigra, and cortex.
GLP-1 receptor agonists (GLP-1RAs; e.g., exenatide, liraglutide,
semaglutide) cross or signal across the blood--brain barrier and have
demonstrated neuroprotective effects in preclinical AD/PD models and in
some clinical trials \citep{hong2024}.

The proposed mechanistic pathway is that GLP-1RAs restore central insulin
sensitivity/signaling (PI3K--Akt pathway, GSK-3$\beta$ inhibition, reduced
neuroinflammation), producing a downstream reduction in pathological protein
aggregation and neuronal loss, and ultimately clinical neuroprotection
(cognitive or motor outcomes). This report evaluates the evidence base for
central insulin resistance as the mediating mechanism, as opposed to GLP-1RA
effects on neuroprotection occurring independently of insulin signaling
(e.g., via direct anti-inflammatory, anti-apoptotic, or mitochondrial
effects).

\subsection{Research question}
What evidence supports central insulin resistance as a mediator of GLP-1
receptor agonist neuroprotection in Alzheimer's and Parkinson's disease?

\subsection{Working hypotheses}
Three candidate hypotheses were formulated and evaluated against the
evidence synthesized in this report (Table~\ref{tab:hypotheses}).

\begin{table}[h]
\centering
\caption{Working hypotheses evaluated in this report.}
\label{tab:hypotheses}
\begin{tabular}{@{}p{0.06\textwidth}p{0.45\textwidth}p{0.4\textwidth}@{}}
\toprule
\textbf{ID} & \textbf{Hypothesis} & \textbf{Summary assessment} \\
\midrule
H1 & GLP-1RAs exert neuroprotective effects \emph{primarily} through
restoration of PI3K--Akt signaling impaired by central insulin resistance. &
Contradicted by decoupling evidence; overclaims a single dominant mechanism. \\[4pt]
H2 & Metabolic dysfunction (diabetes/obesity/insulin resistance) identifies
a responder subgroup for GLP-1 therapy in neurodegenerative disease. &
Weakest support; the one direct empirical test found no such association,
and most trials structurally exclude the relevant population. \\[4pt]
H3 & Brain insulin resistance represents a shared pathological mechanism
linking AD and PD and contributes to (not necessarily primarily explains)
the therapeutic effects of GLP-1RAs. & Best-supported framing; appropriately
scoped, not directly contradicted by available evidence, but not yet
confirmed via formal mediation testing. \\
\bottomrule
\end{tabular}
\end{table}

% =============================================================================
\section{Methods}
\label{sec:methods}

\subsection{Literature search and structured extraction}
% TODO: state exact databases, date range, and search terms used.
Literature was identified via PubMed, PubMed Central, and clinical trial
registries (ClinicalTrials.gov). Search results were supplemented by queries
of the Gene Expression Omnibus (GEO), KEGG pathway database, DrugBank,
OpenAlex, and the Cochrane Library. Eligible study types included randomized
controlled trials, observational/real-world cohort studies, preclinical
(animal/cell) mechanistic studies, and systematic reviews/meta-analyses
reporting either a central insulin-resistance/signaling marker, a
GLP-1RA-related intervention, and/or a neurodegenerative outcome in AD or
PD. For each included source, the following were extracted into a
structured schema: disease studied, GLP-1 agonist used, evidence of insulin
resistance, insulin-signaling markers measured, neurodegenerative outcomes,
whether formal mediation was tested, and main conclusions ($n = 21$ primary
sources; full extraction available in supplementary data).

\subsection{Original transcriptomic re-analysis}
To complement the literature synthesis with original primary-data evidence
on the insulin-resistance side of the hypothesis, we re-analyzed a public
gene-expression dataset, GEO accession GSE34451 (Agilent rat whole-genome
microarray, platform GPL15011), comprising hippocampus, prefrontal cortex,
and striatum samples from control (Wistar), type 1 diabetic
(streptozotocin-treated), and type 2 diabetic (Goto-Kakizaki) rats
($n=3$ per group per region). We focused on the hippocampus, Control vs.\
type 2 diabetes (T2D) contrast as most directly relevant to the
insulin-resistance-in-AD hypothesis.

Microarray probe identifiers were mapped to gene symbols using the GPL15011
platform annotation table, with multi-probe genes collapsed by averaging.
Differential expression between groups was assessed via Welch's $t$-test
with Benjamini--Hochberg false discovery rate (FDR) correction. Genes were
cross-referenced against the KEGG insulin signaling pathway gene set
(map04910/rno04910) to identify insulin-pathway-specific differential
expression. % TODO: cite your own analysis script/repository here, e.g.:
% \citep{yourrepo2026} or a footnote with the GitHub URL.

\subsection{Causal mediation-analysis framework}
No identified study in the literature search performed a formal statistical
test of mediation (treatment $\rightarrow$ insulin-resistance marker
$\rightarrow$ neurodegenerative outcome). We therefore implemented a
Baron--Kenny path-model mediation framework with nonparametric
(case-resampling) bootstrap confidence intervals for the indirect effect,
designed to be applied to individual-patient data the moment such a dataset
becomes available. In the absence of such data, we (a) constructed a
calibrated illustrative simulation using the one available set of real
aggregate statistics in this literature (exenatide-PD trial motor and
biomarker effect sizes), explicitly flagging which parameters were real vs.\
assumed, and (b) conducted a power/feasibility analysis estimating whether
trials of the sample sizes actually used in this literature were capable of
detecting mediation of a plausible effect size.

% =============================================================================
\section{Results}
\label{sec:results}

\subsection{Evidence synthesis from the literature}
Across the 21 primary sources reviewed (7 randomized controlled trials, 7
preclinical animal studies, 4 systematic reviews/meta-analyses including one
Cochrane review, 1 narrative-review-only mechanistic source, and 2
real-world observational cohorts), no source -- clinical or preclinical --
performed a formal statistical mediation analysis (e.g., bootstrapped
indirect-effect testing) of central insulin resistance as a mediator of
GLP-1RA neuroprotective outcomes. Evidence for the mediation hypothesis is
currently mechanistic and correlational rather than causally established.
Mechanistic convergence on shared insulin-signaling components (PI3K/Akt,
GSK-3$\beta$, IRS-1 phosphorylation) was discussed in roughly half of
sources (10/21), concentrated in preclinical studies and narrative reviews.
The clearest direct biomarker evidence in a living human cohort is
\citet{athauda2019}, a secondary exosome-based analysis of the Exenatide-PD
trial reporting increased tyrosine-phosphorylated IRS-1, total Akt, and
phosphorylated mTOR with exenatide, temporally associated with motor
improvement -- but without formal mediation testing. Clinical trial primary
endpoints were directionally inconsistent by disease: all three Parkinson's
disease RCTs identified (exenatide \citep{avilesolmos2013,athauda2017} and
lixisenatide \citep{meissner2024}) were positive on the primary motor
endpoint, whereas the three Alzheimer's disease RCTs reviewed
(\citet{gejl2016}, \citet{edison2025}/ELAD, \citet{cummings2026}/evoke and
evoke+) were null or mixed on their primary endpoints. Pooled
meta-analytic clinical evidence was null or low-certainty in both diseases
(\citet{santos2026}: SMD $-0.21$, ns, for AD/MCI cognitive outcomes;
\citet{helal2025}: only 1 of 6 pooled PD outcomes significant;
\citet{mulvaney2020}, Cochrane: GRADE certainty low-to-very-low).

\begin{table}[h]
\centering
\caption{Approximate tally of evidence type and direction across the 21
reviewed sources. Categories are not mutually exclusive -- a single source
may contribute to more than one row.}
\label{tab:evidence-summary}
\begin{tabular}{@{}lcc@{}}
\toprule
\textbf{Evidence category} & \textbf{Supporting} & \textbf{Contradicting/null} \\
\midrule
Mechanistic convergence (PI3K/Akt, GSK-3$\beta$ overlap) & 10 & -- \\
Preclinical marker--phenotype co-movement & 2 & 2 \\
Clinical biomarker evidence (human, in vivo) & 1 & -- \\
Primary clinical trial endpoints (PD positive / AD null-mixed) & 3 & 3 \\
Pooled meta-analytic clinical evidence & -- & 3 \\
Metabolic-subgroup (responder) analyses & -- & 1 \\
\bottomrule
\end{tabular}
\end{table}

\subsection{Contradicting and decoupling evidence}
Seven distinct ways in which the available evidence complicates or
contradicts a simple mediation chain were identified (full detail in
\texttt{extraction/contradicting\_evidence.json}). The single cleanest
within-study dissociation is \citet{carranza2021}: liraglutide produced the
full expected neuroprotective phenotype (reduced cortical atrophy, amyloid
burden, tau hyperphosphorylation, microglial activation, and improved
cognition) in a mixed AD/T2D mouse model, yet cortical IR-A, IR-B, and
IGF-1R mRNA expression showed \emph{no difference} between groups -- the
proposed mediator did not move even though the proposed downstream effect
did. \citet{paladugu2021} showed the converse partial dissociation: liraglutide
restored IDE and phosphorylated GSK-3$\beta$ but explicitly did not elevate
phosphorylated Akt, and neither degenerating-neuron counts nor behavioral
measures changed. At the clinical level, the trials with the most literal
central-insulin/glucose-metabolism endpoint (cerebral glucose metabolism via
PET) were null on that specific endpoint in both \citet{gejl2016} and
\citet{edison2025} even as other benefits (ADAS-Exec, brain volume) appeared
through a different readout, and the largest, most metabolically diverse AD
trial (\citet{cummings2026}, evoke/evoke+, $n=3{,}808$) was completely null
on its primary endpoint. The most direct empirical test of patient-level
prediction, the \citet{athauda2019} post-hoc analysis of the Exenatide-PD
cohort, examined BMI and insulin resistance as candidate predictors of motor
response and found neither significant; tremor-dominant phenotype and
milder baseline disease predicted benefit instead.

\subsection{Metabolic-subgroup analysis}
Five of the seven RCTs identified (\citet{athauda2017}, Aviles-Olmos 2013,
\citet{gejl2016}, \citet{edison2025}, and very likely \citet{meissner2024})
explicitly excluded diabetic participants at enrollment, specifically to
isolate a CNS drug effect from peripheral glycemic confounding -- which
means most of this evidence base structurally cannot test whether
metabolically impaired patients benefit more, because that subgroup was
screened out before randomization. Only \citet{cummings2026} (evoke and
evoke+) enrolled a non-trivial diabetic subgroup (13.6\% overall, $n
\approx 516$), but its primary outcome was null overall and no
diabetes-stratified subgroup analysis has been published in any accessible
source. The one direct test of metabolic status as a predictor of GLP-1RA
response (\citet{athauda2019} post-hoc analysis) found that BMI and insulin
resistance were not significant predictors of motor benefit in PD. The
largest dataset with substantial metabolically-impaired representation
(AbuAlrob et al.\ 2025, real-world cohort, $n=295{,}010$, $\sim$51\%
diabetic in the treated arm) had the statistical power to test this
directly but did not perform the stratified analysis, citing incomplete EHR
coding of diabetes/agent/dose details. The honest state of the evidence is
therefore an evidence gap, not a documented null or positive subgroup
effect: trial designs have systematically avoided enrolling the population
whose response would test the metabolic-subgroup hypothesis (H2) most
directly.

\subsection{Original transcriptomic analysis (GSE34451)}

\begin{figure}[h]
\centering
\includegraphics[width=0.92\textwidth]{Differential expression.png}
\caption{Volcano plot of genome-wide differential expression between control
and type 2 diabetic (Goto-Kakizaki) rat hippocampus (GSE34451). KEGG
insulin-signaling pathway genes (rno04910, $n=125$ annotated in this
dataset) highlighted in red; none reach the FDR-significance line at this
sample size. $n=3$ per group.}
\label{fig:volcano}
\end{figure}

\begin{figure}[h]
\centering
\includegraphics[width=0.92\textwidth]{Top 20 insulin signalling.png}
\caption{Per-replicate expression heatmap of the top 20 insulin-signaling
pathway genes ranked by raw (uncorrected) $p$-value, mean-centered per gene.
Note the substantial within-group heterogeneity across replicates (e.g.\
T2D replicates disagree in direction for several genes), consistent with
the small per-group sample size.}
\label{fig:heatmap}
\end{figure}

\begin{figure}[h]
\centering
\includegraphics[width=0.92\textwidth]{Box plots.png}
\caption{Expression distributions for the six most nominally significant
insulin-signaling-pathway genes (raw $p<0.06$) by group: \textit{Insr},
\textit{Pde3b}, \textit{Gys2}, \textit{Prkab2}, \textit{Mapk8},
\textit{Ppargc1a}.}
\label{fig:boxplots}
\end{figure}

\begin{figure}[h]
\centering
\includegraphics[width=0.7\textwidth]{PCA.png}
\caption{Principal component analysis of hippocampus samples, colored by
group. PC1 (42.2\% of variance) shows a partial trend separating Control
from T2D samples, with some overlap, consistent with a modest but
incomplete global transcriptomic difference at this sample size.}
\label{fig:pca}
\end{figure}

\begin{figure}[h]
\centering
\begin{subfigure}{0.48\textwidth}
  \includegraphics[width=\textwidth]{p-value distribution.png}
  \caption{Genome-wide raw $p$-value distribution}
\end{subfigure}
\hfill
\begin{subfigure}{0.48\textwidth}
  \includegraphics[width=\textwidth]{Sample correlation.png}
  \caption{Sample--sample correlation matrix}
\end{subfigure}
\caption{(a) The raw $p$-value histogram shows a modest excess of small
$p$-values above what a uniform null distribution would predict, indicating
weak but non-zero genome-wide signal that is underpowered for FDR
significance at $n=3$/group. (b) Sample-sample Pearson correlations are
uniformly high (0.99--1.00) with no sharp block structure separating
Control from T2D, typical of normalized microarray data and consistent with
the modest global separation seen in the PCA.}
\label{fig:pval-corr}
\end{figure}

\begin{figure}[h]
\centering
\includegraphics[width=0.92\textwidth]{Forest plots.png}
\caption{Effect sizes (log$_2$ fold change, T2D vs.\ Control) with
confidence intervals across the annotated insulin-signaling gene panel,
ranked by raw $p$-value; orange bars indicate raw $p<0.05$.}
\label{fig:forest}
\end{figure}

Of 125 KEGG insulin-signaling pathway genes (rno04910) annotated in this
dataset, none survived Benjamini--Hochberg FDR correction, but five reached
nominal significance (raw $p<0.05$): \textit{Insr} (log$_2$FC $=+0.76$,
raw $p=0.0067$, the single most significant hit in the entire dataset),
\textit{Pde3b} (log$_2$FC $\approx -1.3$, raw $p=0.0141$), \textit{Gys2}
(log$_2$FC $\approx +0.37$, raw $p=0.0275$), \textit{Prkab2} (log$_2$FC
$\approx -0.95$, raw $p=0.0299$), and \textit{Mapk8}/JNK1 (log$_2$FC
$\approx +0.36$, raw $p=0.0493$); \textit{Ppargc1a} narrowly missed
significance (raw $p=0.0536$). Direction of effect was mixed relative to
the canonical insulin-resistance model: \textit{Mapk8} upregulation in the
diabetic group is consistent with JNK-mediated inhibitory serine
phosphorylation of IRS-1, a well-established driver of insulin resistance,
and \textit{Pde3b}/\textit{Prkab2} downregulation is consistent with
reduced downstream Akt/AMPK pathway output. By contrast, \textit{Insr} and
\textit{Gys2} were both \emph{upregulated} in the diabetic group, which is
not the naively expected direction (receptor and glycogen-synthase
transcripts are often assumed to decrease alongside impaired insulin
action). \textit{Insr} upregulation has a plausible compensatory
explanation -- receptor transcript upregulation as homeostatic feedback when
downstream signaling efficiency drops -- but this was not directly tested
in the present dataset.

\subsection{Mediation feasibility and power analysis}

\begin{figure}[h]
\centering
\includegraphics[width=0.92\textwidth]{illustrative.png}
\caption{Illustrative mediation path model calibrated to the aggregate
statistics reported in the Exenatide-PD trial and its exosome sub-study.
Path~$a$ (IRS-1 phosphorylation change, $+0.279$, $p=0.000106$) and the
total/direct effect ($c' = +5.15$ MDS-UPDRS III points) are real published
aggregate statistics; path~$b$ ($+3.99$, $p=0.109$) is an assumed
illustrative value, since no published source reports a per-patient
biomarker--outcome regression coefficient (see Methods). The resulting
illustrative indirect effect ($a \times b = +1.11$, 95\% CI $[0.0593,
2.34]$, excluding zero) is a feasibility demonstration of the analysis
framework, \emph{not} a real mediation finding.}
\label{fig:mediation-illustration}
\end{figure}

\begin{figure}[h]
\centering
\includegraphics[width=0.85\textwidth]{mediation_power_curve.png}
\caption{Empirical statistical power to detect a hypothetical mediation
effect (bootstrapped indirect-effect 95\% CI excluding zero) as a function
of trial sample size, at assumed parameters $a=0.27$, $b=0.5$, $c'=1.75$
(see Methods/Section~\ref{fig:mediation-illustration}). Evaluated at sample
sizes spanning those actually used in this literature ($n=62$, the
Exenatide-PD trial size, through $n=3{,}808$, the evoke/evoke+ scale).
Dashed line marks the conventional 80\% power threshold.}
\label{fig:power-curve}
\end{figure}

% NOTE: values below are read off the plotted power curve (approximate);
% replace with exact decimals from output/power/power_analysis.csv if/when
% available for full precision.
At the assumed effect size, the Exenatide-PD trial scale ($n=62$) had only
$\approx$14\% power to detect mediation via this framework -- meaning even
if a mediated effect of this magnitude were truly present in that trial, it
would very likely have gone undetected by a formal mediation test. Power
rose to $\approx$70\% at $n=300$ and reached $\approx$100\% by $n=1{,}200$,
remaining effectively at ceiling at the evoke/evoke+ scale ($n=3{,}808$);
the conventional 80\% power threshold was crossed somewhere between
$n=300$ and $n=1{,}200$. This indicates that, at this assumed effect size,
only the largest trials conducted in this literature (evoke/evoke+, and by
extension ELAD at $n=204$ only marginally) were ever statistically
positioned to detect a mediated effect with reasonable power -- the smaller
proof-of-concept PD trials (Exenatide-PD, $n=62$; LixiPark, $n\approx156$)
were comparable in scale and thus a true mediation effect of this size could
easily have escaped statistical detection in them even if biologically
present.

% =============================================================================
\section{Discussion}
\label{sec:discussion}

\subsection{Which hypothesis is best supported}
Of the three candidate hypotheses, H3 -- that central insulin resistance is
a shared, contributory (rather than primary or subgroup-defining) mechanism
linking AD and PD to GLP-1RA effects -- is the framing best supported by,
and least contradicted by, the available evidence. H1 (insulin-resistance
reversal as the \emph{primary} mechanism) is directly contradicted by the
decoupling evidence in Section~\ref{sec:results}: \citet{carranza2021}
showed a full neuroprotective phenotype with a completely flat
insulin-receptor mRNA readout, and \citet{paladugu2021} showed biochemical
pathway engagement (pGSK-3$\beta$, IDE) without a corresponding behavioral or
neurodegenerative-marker change in the same animals -- if insulin-resistance
reversal were the dominant mechanism, these dissociations should not occur
as readily as they do. H2 (metabolic status as a subgroup-defining
predictor) is unsupported by the one direct empirical test available: the
\citet{athauda2019} post-hoc analysis found neither BMI nor insulin
resistance among the significant predictors of motor response, with
phenotype and disease-stage variables predicting response instead. H3
survives both of these challenges because it does not require insulin-pathway
engagement to be necessary or sufficient in every case -- it only requires
that the pathway be \emph{one real contributor among several}, which is
consistent with cases where insulin-signaling biomarkers and clinical
benefit do move together (\citet{athauda2019} exosome substudy; Zhang et
al.\ 2022, 6-OHDA rat model, pIRS-1$^{\text{Ser312}}$ normalization
co-occurring with dopaminergic neuron preservation) as well as cases where
they dissociate.

\subsection{Mechanistic plausibility versus clinical translation}
Preclinical mechanistic convergence on insulin-signaling components is
broad (10/21 sources) and preclinical pooled effect sizes are large
(\citet{kong2023}: SMD $-1.28$ to $-1.88$ across learning, memory, and
amyloid/tau outcomes), yet clinical translation is markedly inconsistent
between the two diseases: all three Parkinson's disease RCTs identified
were positive on the primary motor endpoint, while all three Alzheimer's
disease RCTs reviewed were null or mixed on their primary endpoints. Two
non-mutually-exclusive explanations are consistent with the data. First, no
clinical trial in either disease actually measured central insulin
resistance or signaling as an outcome (with the partial exception of the
exosome substudy in PD), so a clinical translation gap cannot currently be
attributed to mechanism failure versus simply absent measurement -- the
trials were not designed to test the mechanism, only the downstream
clinical endpoint. Second, disease-specific factors plausibly modulate
translation independently of insulin signaling: PD's relatively focal
dopaminergic motor circuitry may be more tractable to a partial, time-limited
biochemical intervention than AD's diffuse cortical pathology, and
blood--brain-barrier penetration differences between agents (e.g.\ DA5-CH's
superiority over liraglutide in \citet{zhang2023}, attributed by the authors
to penetration rather than differential insulin-pathway engagement) suggest
that drug-delivery pharmacokinetics, not insulin-pathway engagement per se,
may be the dominant driver of between-agent differences observed to date.

\subsection{The unanswered metabolic-subgroup question}
Most trials cannot test H2 because diabetic and severely insulin-resistant
participants were excluded from enrollment by design in 5 of 7 RCTs
identified, specifically to isolate a CNS drug effect from peripheral
glycemic confounding. This is a reasonable design choice for isolating a
CNS-specific signal, but it has the side effect of removing exactly the
population in which a true insulin-resistance-mediated effect should be
largest and easiest to detect. Two paths exist to close this gap with data
that may already exist: (1) a post-hoc, diabetes/HbA1c-stratified
reanalysis of \citet{cummings2026} (evoke/evoke+), which enrolled
$\sim$516 diabetic participants -- the only GLP-1RA AD/PD trial with a
metabolically-impaired subgroup of meaningful size -- but has not published
such a stratified result; and (2) a diabetes-stratified reanalysis of the
AbuAlrob et al.\ 2025 real-world cohort ($n=295{,}010$, $\sim$51\% diabetic
in the treated arm), which the original authors explicitly did not perform
due to incomplete EHR coding. Absent either, a prospective trial enrolling
a broader metabolic range with a central insulin-resistance biomarker (e.g.\
CSF or neuronal-exosome IRS-1/Akt/mTOR markers, as in \citet{athauda2019})
measured at baseline and follow-up would directly test H2 rather than
inferring it indirectly.

\subsection{Interpretation of the original transcriptomic finding}
The GSE34451 re-analysis is a genuine, original contribution addressing the
insulin-resistance side of the hypothesis in a relevant tissue (hippocampus)
and disease-relevant model (Goto-Kakizaki type 2 diabetic rat), but it has
an important scope limitation: this dataset contains no GLP-1RA treatment
arm, so it cannot itself test mediation -- it can only establish whether
insulin-pathway transcriptional changes are present in a diabetic
hippocampus, which is the upstream half of the proposed causal chain. Within
that scope, the finding is biologically plausible but mixed. \textit{Mapk8}
upregulation matches the canonical JNK-mediated insulin-resistance
mechanism described in the broader literature \citep{talbot2014}, lending
some transcriptomic support to the idea that this pathway is engaged in a
diabetic hippocampus. However, \textit{Insr} and \textit{Gys2} upregulation
runs counter to the simplest version of the hypothesis (in which insulin
signaling components would be expected to decrease uniformly). A plausible
resolution is that mRNA abundance does not cleanly track pathway
\emph{activity}: \citet{talbot2014} reviews postmortem human hippocampal
tissue-stimulation assays showing reduced insulin-stimulated IR and IRS-1
\emph{tyrosine phosphorylation} in AD relative to baseline transcript or
protein levels, i.e.\ the functional deficit in this literature is
predominantly post-translational (phosphorylation-state), not
transcriptional. \textit{Insr} mRNA upregulation alongside reduced
receptor-level signaling efficiency would be consistent with a compensatory
feedback loop rather than a contradiction, though this was not directly
tested here. These findings should be treated as hypothesis-generating: the
$n=3$/group sample size limits statistical power substantially, and zero
genes survived FDR correction at the genome-wide level.

% =============================================================================
\section{Limitations}
\label{sec:limitations}

\begin{itemize}[leftmargin=*]
  \item No identified study -- across 21 primary sources spanning randomized
  trials, preclinical mechanistic studies, meta-analyses, and a Cochrane
  review -- performed a formal statistical mediation analysis linking
  GLP-1RA treatment, a central insulin-resistance marker, and a
  neurodegenerative outcome in the same subjects.
  \item No public dataset (literature- or GEO-derived) links GLP-1RA
  treatment, a central insulin-resistance marker, and a neurodegenerative
  outcome in the same subjects; the original transcriptomic analysis in this
  report addresses only the insulin-resistance side of the hypothesis.
  \item Most randomized trials identified excluded diabetic participants by
  design, structurally preventing direct tests of the metabolic-subgroup
  hypothesis (H2).
  \item The original transcriptomic analysis (GSE34451) used $n=3$ per
  group per brain region, limiting statistical power; zero genes survived
  genome-wide FDR correction, and findings should be treated as
  exploratory/hypothesis-generating rather than confirmatory.
  \item The nominally significant insulin-pathway genes identified
  (\textit{Insr}, \textit{Pde3b}, \textit{Gys2}, \textit{Prkab2},
  \textit{Mapk8}) did not all change in the direction predicted by the
  simplest version of the hypothesis (\textit{Insr} and \textit{Gys2} were
  upregulated, not downregulated, in the diabetic group), so this result
  should not be over-interpreted as clean confirmatory evidence even within
  its own limited scope.
  \item GSE34451 contains no GLP-1RA treatment arm and therefore cannot
  itself test mediation; it addresses only the insulin-resistance side of
  the proposed causal chain.
  \item KEGG's curated reference pathways for AD and PD do not currently
  integrate insulin-signaling components, indicating this mechanism is not
  yet reflected in canonical pathway curation.
  \item The power-analysis values reported in Section~\ref{sec:results} are
  read directly off the plotted power curve rather than the underlying
  \texttt{power\_analysis.csv} decimals, and the curve itself is a
  feasibility calculation under assumed (not directly observed) path
  coefficients -- it estimates detectability of a hypothesized effect size,
  not whether that effect size is real.
\end{itemize}

% =============================================================================
\section{Conclusion}
\label{sec:conclusion}
Formal mediation evidence linking GLP-1RA treatment, central insulin
resistance, and neurodegenerative outcomes does not yet exist anywhere in
the reviewed literature: no clinical or preclinical study has performed a
statistical mediation test of this causal chain. Among the three candidate
hypotheses evaluated, H3 -- central insulin resistance as a shared,
contributory (not primary or subgroup-defining) mechanism -- is the
best-supported framing, because it is the only one consistent with both the
cases where insulin-pathway engagement and clinical benefit co-occur and the
cases where they clearly dissociate. Our original re-analysis of GSE34451
adds a modest, hypothesis-generating data point on the insulin-resistance
side of this chain: insulin-signaling-pathway genes show nominal (though not
FDR-significant) differential expression in diabetic rat hippocampus, with
a mixed direction of effect that itself argues against an oversimplified,
uniformly-down-regulated picture of central insulin resistance. Closing the
remaining evidence gap requires a dataset or sub-study that measures GLP-1RA
exposure, a central insulin-resistance biomarker, and a neurodegenerative
outcome in the same subjects -- the most immediately actionable opportunity
being a diabetes/HbA1c-stratified reanalysis of the evoke/evoke+ trial
cohort, which already contains the necessary metabolic subgroup but has not
published a stratified result.

% =============================================================================
\section*{Acknowledgments}
% TODO

\section*{Data and code availability}
% TODO: insert your GitHub repository URL.
All extraction files, analysis code, and figure-generation scripts are
available at: \url{https://github.com/Swdvvv/Insulin-resistance-as-a-mediator-of-GLP-1-agonist-}

% =============================================================================
\bibliographystyle{plainnat}
\begin{thebibliography}{99}

\bibitem{athauda2017} Athauda D, Maclagan K, Skene SS, et al. Exenatide once
weekly versus placebo in Parkinson's disease: a randomised, double-blind,
placebo-controlled trial. \textit{Lancet}. 2017;390(10103):1664--1675.

\bibitem{athauda2019} Athauda D, Gulyani S, Karnati HK, et al. Utility of
Neuronal-Derived Exosomes to Examine Molecular Mechanisms That Affect Motor
Function in Patients With Parkinson Disease: A Secondary Analysis of the
Exenatide-PD Trial. \textit{JAMA Neurol}. 2019;76(4):420--429.

\bibitem{avilesolmos2013} Aviles-Olmos I, Dickson J, Kefalopoulou Z, et al.
Exenatide and the treatment of patients with Parkinson's disease.
\textit{J Clin Invest}. 2013;123(6):2730--2736.

\bibitem{meissner2024} Meissner WG, Remy P, Giordana C, et al. Trial of
Lixisenatide in Early Parkinson's Disease. \textit{N Engl J Med}.
2024;390(13):1176--1185.

\bibitem{gejl2016} Gejl M, Gjedde A, Egefjord L, et al. In Alzheimer's
Disease, 6-Month Treatment with GLP-1 Analog Prevents Decline of Brain
Glucose Metabolism: Randomized, Placebo-Controlled, Double-Blind Clinical
Trial. \textit{Front Aging Neurosci}. 2016;8:108.

\bibitem{edison2025} Edison P, Femminella GD, Ritchie C, et al. Liraglutide
in mild to moderate Alzheimer's disease: a phase 2b clinical trial.
\textit{Nat Med}. 2025;32(1):353--361.

\bibitem{cummings2026} Cummings JL, Atri A, Sano M, et al. Efficacy and
safety of oral semaglutide 14~mg (flexible dose) in early-stage symptomatic
Alzheimer's disease (evoke and evoke+): two phase 3, randomised,
placebo-controlled trials. \textit{Lancet}. 2026;407(10544):2167--2179.

\bibitem{hansen2015} Hansen HH, et al. The GLP-1 Receptor Agonist
Liraglutide Improves Memory Function and Increases Hippocampal CA1
Neuronal Numbers in a Senescence-Accelerated Mouse Model of Alzheimer's
Disease. \textit{J Alzheimers Dis}. 2015;46(4):877--888.

\bibitem{batista2018} Batista AF, Forny-Germano L, Clarke JR, et al. The
diabetes drug liraglutide reverses cognitive impairment in mice and
attenuates insulin receptor and synaptic pathology in a non-human primate
model of Alzheimer's disease. \textit{J Pathol}. 2018;245(1):85--100.

\bibitem{paladugu2021} Paladugu L, Gharaibeh A, Kolli N, et al. Liraglutide
Has Anti-Inflammatory and Anti-Amyloid Properties in Streptozotocin-Induced
and 5xFAD Mouse Models of Alzheimer's Disease. \textit{Int J Mol Sci}.
2021;22(2):860.

\bibitem{mcclean2014} McClean PL, H\"olscher C. Liraglutide can reverse
memory impairment, synaptic loss and reduce plaque load in aged APP/PS1
mice, a model of Alzheimer's disease. \textit{Neuropharmacology}.
2014;76(Pt A):57--67.

\bibitem{carranza2021} Carranza-Naval MJ, del Marco A, Hierro-Bujalance C,
et al. Liraglutide reduces vascular damage, neuronal loss, and cognitive
impairment in a mixed murine model of Alzheimer's disease and type 2
diabetes. \textit{Front Aging Neurosci}. 2021;13:741923.

\bibitem{zhang2023} Zhang Z, et al. A Dual GLP-1/GIP Receptor Agonist Is
More Effective than Liraglutide in the A53T Mouse Model of Parkinson's
Disease. \textit{Parkinsons Dis}. 2023;2023:7427136.

\bibitem{zhang2022} Zhang L, et al. DA5-CH and Semaglutide Protect against
Neurodegeneration and Reduce $\alpha$-Synuclein Levels in the 6-OHDA
Parkinson's Disease Rat Model. \textit{Parkinsons Dis}. 2022;2022:1428817.

\bibitem{talbot2014} Talbot K. Brain insulin resistance in Alzheimer's
disease and its potential treatment with GLP-1 analogs. \textit{Neurodegener
Dis Manag}. 2014;4(1):31--40.

\bibitem{kong2023} Kong F, et al. Glucagon-like peptide 1 (GLP-1) receptor
agonists in experimental Alzheimer's disease models: a systematic review and
meta-analysis of preclinical studies. \textit{Front Pharmacol}.
2023;14:1205207.

\bibitem{hong2024} Hong CT, Chen JH, Hu CJ. Role of glucagon-like
peptide-1 receptor agonists in Alzheimer's disease and Parkinson's disease.
\textit{J Biomed Sci}. 2024;31:102.

\bibitem{helal2025} Helal MM, AbouShawareb H, Abbas OH, et al. GLP-1
receptor agonists in Parkinson's disease: an updated comprehensive
systematic review with meta-analysis. \textit{Diabetol Metab Syndr}.
2025;17:352.

\bibitem{mulvaney2020} Mulvaney CA, Duarte GS, Handley J, et al. GLP-1
receptor agonists for Parkinson's disease. \textit{Cochrane Database Syst
Rev}. 2020;7(7):CD012990.

\bibitem{santos2026} Santos P, S\'a Filho AS, Aprigliano V, et al.
Liraglutide and Exenatide in Alzheimer's Disease and Mild Cognitive
Impairment: A Systematic Review and Meta-Analysis of Cognitive Outcomes.
\textit{Pharmaceutics}. 2026;18(1):69.

\bibitem{abualrob2025} AbuAlrob MA, Itbaisha A, Abujwaid YK, et al.
Exploring the neuroprotective role of GLP-1 agonists against Alzheimer's
disease: Real-world evidence from a propensity-matched cohort. \textit{J
Alzheimers Dis Rep}. 2025;9:25424823251388650.

\bibitem{wang2025} Wang W, Davis PB, Qi X, et al. Associations of
semaglutide with Alzheimer's disease-related dementias in patients with
type 2 diabetes: A real-world target trial emulation study. \textit{J
Alzheimers Dis}. 2025;106(4):1509--1522.

\end{thebibliography}

\end{document}
